\documentclass[11pt,a4paper]{article}
\usepackage[T1]{fontenc}
\usepackage{lmodern}
\usepackage[margin=26mm]{geometry}
\usepackage{amsmath,amssymb,amsthm,booktabs,array,longtable}
\usepackage{microtype,multicol}
\usepackage[hidelinks]{hyperref}
\setlength{\parskip}{3pt}
\setlength{\emergencystretch}{2em}
\raggedbottom
\allowdisplaybreaks[1]
\clubpenalty=10000
\widowpenalty=10000
\newtheorem{theorem}{Theorem}[section]
\newtheorem{lemma}[theorem]{Lemma}
\newtheorem{corollary}[theorem]{Corollary}
\theoremstyle{remark}
\newtheorem{remark}[theorem]{Remark}
\newcommand{\Z}{\mathbb Z}
\newcommand{\Q}{\mathbb Q}
\newcommand{\N}{\mathbb N}
\newcommand{\F}{\mathbb F}
\newcommand{\ps}[1]{#1[[x]]}
\newcommand{\Rev}{\operatorname{Rev}}
\newcommand{\rev}{^{\langle-1\rangle}}
\newcommand{\Plus}{\operatorname{Plus}}
\newcommand{\Minus}{\operatorname{Minus}}
\newcommand{\coeff}[2]{[x^{#1}]#2}
\newcommand{\E}{\mathcal E}
\Urlmuskip=0mu plus 1mu
\hypersetup{pdftitle={A389542: A reversion that becomes an involution modulo four},pdfauthor={Dev Goyal},pdfsubject={Separate manuscript with AI disclosure and Lean proof map}}
\title{A389542\\[4pt]\large A reversion that becomes an involution modulo four}
\author{Dev Goyal}
\date{24 October 2026}
\begin{document}
\maketitle
\begin{center}
\fbox{\begin{minipage}{0.9\linewidth}\small
\textbf{AI disclosure.} ChatGPT (OpenAI) generated the mathematical arguments,
Lean proofs, and this manuscript in response to Dev Goyal's requests.
Dev Goyal is the sole human author. The source sequence and conjecture or
observation are credited to Paul D. Hanna.
\textbf{Independent human review is pending.}
\end{minipage}}
\end{center}
\begin{abstract}
We construct the unique normalized integer solution of $A(A^2-x^2)=4A^3$. Above degree one, its residue modulo four is two exactly at indices $2^k+1$ and zero otherwise. An auxiliary inverse becomes an involution after reduction; no cancellation of two in a residue ring is used.
\end{abstract}

\section{Source, notation, and scope}

The source is Paul D. Hanna's OEIS entry A389542 \cite{oeis542},
retrieved on 24 October 2026. Its initial sequence index is 1.
We also retain the formal zero constant coefficient at index zero.

All series are formal. The set $\N$ includes zero, and $[x^n]U$ denotes a
coefficient. A normalized series is $x+O(x^2)$; $O(x^r)$ means divisibility
by $x^r$, over any coefficient ring used below. We write $U\circ V=U(V(x))$
only when $V(0)=0$. Compositional inversion, denoted $\Rev(U)$ or $U\rev$,
is distinct from multiplicative inversion of a unit. Congruences are
coefficientwise; ordinary products and compositional iterates are distinct.

\section{Formal algebra used in the proof}

\begin{lemma}[Coefficientwise construction]\label{lem:contraction}
Let $K$ be a commutative ring. Suppose an operator $\Phi$ preserves
normalized series and, for $r\ge2$, satisfies
\[
 G-H=O(x^r)\quad\Longrightarrow\quad
 \Phi(G)-\Phi(H)=O(x^{r+1}).
\]
Then $\Phi$ has exactly one normalized fixed point in $\ps K$.
More generally, a residual equation whose degree-$r$ condition determines
the next unknown coefficient with multiplier one has a unique solution
once its initial coefficients are fixed and its lower-degree conditions
hold identically.
\end{lemma}

\begin{proof}
Begin with $G_0=x$ and set $G_{m+1}=\Phi(G_m)$. Normalization gives
$G_1-G_0=O(x^2)$. Inductively,
$G_{m+1}-G_m=O(x^{m+2})$, so each coefficient eventually stabilizes.
Define $G$ using those stable coefficients. The displayed implication
ensures that each coefficient of $\Phi(G)$ equals the corresponding stable
coefficient, hence $G=\Phi(G)$. If two fixed points agree below degree
$r$, they agree through degree $r$; induction proves uniqueness.

For the residual version, set the next coefficient to the negative of the
known residual term. Its multiplier is one, so the construction uses only
ring operations. No earlier equation changes, and every equation is
eventually imposed. The same coefficient comparison proves uniqueness.
\end{proof}

\begin{lemma}[Units and reversion]\label{lem:inverse}
Every constant-one series has a unique constant-one multiplicative
inverse. Every normalized series has a unique normalized two-sided
compositional inverse. These inversions, formal differentiation, and
valid substitution commute with coefficient ring homomorphisms.
\end{lemma}

\begin{proof}
For $UV=1$ with $U(0)=1$, the coefficient of degree $n>0$ is $v_n$ plus
a finite expression in earlier coefficients of $V$, which determines
$v_n$ without division. For normalized $G$, the coefficient of degree
$n>1$ in $G(H)$ is $h_n$ plus an expression in earlier coefficients of
$H$, giving a unique normalized right inverse $H$.

To see the second inverse identity, construct a normalized right inverse
$J$ of $H$. Associativity gives
$G=G\circ(H\circ J)=(G\circ H)\circ J=J$, and hence $H\circ G=x$.
The same association proves uniqueness. Each coefficient of a valid
composition is a finite polynomial in coefficients. Differentiation
multiplies the degree-$(n+1)$ coefficient by $n+1$. Ring homomorphisms
therefore preserve these operations and equations; uniqueness identifies
the mapped inverses with the inverses of the mapped series.
\end{proof}

\begin{lemma}[Binary support]\label{lem:binary}
There is a unique normalized integral series $C$ satisfying $C=x+C^2$.
Its reduction modulo two is
\[
 P(x)=\sum_{k\ge0}x^{2^k}.
\]
More generally, for every integer $d\ge1$, the unique zero-constant
solution of $V+V^2=x^d$ in $\ps{\F_2}$ is
\[
 V=P(x^d)=\sum_{k\ge0}x^{d2^k}.
\]
\end{lemma}

\begin{proof}
The operator $G\mapsto x+G^2$ satisfies Lemma~\ref{lem:contraction},
since $G^2-H^2=(G-H)(G+H)$ and $G+H=O(x)$. In characteristic two,
every series $V$ satisfies $V(x)^2=V(x^2)$. This follows by truncating
above any fixed degree and applying the polynomial Frobenius identity.

Thus $V+V^2=x^d$ is a recursion in which the coefficient at an odd positive
index is fixed immediately, while the coefficient at an even index $n$
depends on the coefficient at $n/2<n$. The constant coefficient is fixed
to zero. This proves uniqueness. The proposed sum is well-defined
coefficientwise because its exponents tend to infinity; adding its square
cancels all terms except $x^d$. Set $d=1$ to identify the reduction of $C$.
\end{proof}

\begin{lemma}[Doubling and substitution]\label{lem:double}
For $U,V\in(\Z/4\Z)[[x]]$,
\[
 2U=2V\quad\Longleftrightarrow\quad
 U\equiv V\pmod2.
\]
If $B,D$ have zero constant coefficient and $B\equiv D\pmod2$, then
$2H(B)=2H(D)$ for every series $H$ over $\Z/4\Z$.
\end{lemma}

\begin{proof}
For residues $u,v$ modulo four, $2u=2v$ exactly when their reductions modulo
two agree. Apply this fact to each coefficient. Reduction modulo two
commutes with substitution, giving the second assertion. In particular,
this lemma never cancels two inside $\Z/4\Z$.
\end{proof}

\section{A389542: an integral inverse that becomes an involution}

\subsection{An integral construction and exact source correspondence}

Construct a normalized series $L$ over any commutative ring from
\begin{equation}\label{eq:squareL}
 L=x+L^2-2xL(4x^3).
\end{equation}
The right side preserves normalization. If $G-H=O(x^r)$, then
$G^2-H^2=O(x^{r+1})$ for normalized inputs, while
$x(G-H)(4x^3)=O(x^{3r+1})$. Lemma~\ref{lem:contraction} proves
existence and uniqueness without division. More explicitly, for $n\ge2$,
\[
 \ell_n=\sum_{i=1}^{n-1}\ell_i\ell_{n-i}
 -\begin{cases}
  2\,4^{(n-1)/3}\ell_{(n-1)/3},&3\mid n-1,\\
  0,&3\nmid n-1,
 \end{cases}
 \qquad \ell_1=1.
\]
Set $R=x-2xL$. Equation~\eqref{eq:squareL} gives
\begin{align}
 R^2+R(4x^3)
 &=x^2+4x^2\bigl(-L+L^2+x-2xL(4x^3)\bigr)\notag\\
 &=x^2.\label{eq:squareR}
\end{align}
Let $A=R\rev$, which exists integrally by Lemma~\ref{lem:inverse}.
Composing \eqref{eq:squareR} on the right with $A$ gives
$x^2+R(4A^3)=A^2$. Applying $A$ to $R(4A^3)=A^2-x^2$ yields
the exact source equation
\begin{equation}\label{eq:squaresource}
 A\bigl(A(x)^2-x^2\bigr)=4A(x)^3.
\end{equation}

For completeness, this normalized integer source solution is unique.
If a normalized $U\in\ps{\Z}$ satisfies \eqref{eq:squaresource} and
$V=U\rev$, substitution of $V(t)$ gives
$U(t^2-V(t)^2)=4t^3$. Applying $V$ yields
\begin{equation}\label{eq:squareinverseunique}
 V(t)^2+V(4t^3)=t^2.
\end{equation}
If two normalized integer solutions of \eqref{eq:squareinverseunique}
first differ by $c x^r+O(x^{r+1})$, $r\ge2$, their squares first differ
by $2c x^{r+1}+O(x^{r+2})$. Their substituted difference has order at
least $3r\ge r+2$. Subtracting the equations forces $2c=0$ in $\Z$,
and hence $c=0$, a contradiction. Thus $V=R$ and $U=A$.
The uniqueness argument here is specifically over $\Z$; cancellation of
two is not asserted over rings with two-torsion.

\subsection{Reduction modulo four}

\begin{theorem}\label{thm:squareresidue}
The unique normalized integer solution of \eqref{eq:squaresource} satisfies,
for every $n>1$,
\[
 [x^n]A\equiv
 \begin{cases}
  2\pmod4,&n=2^k+1\text{ for some }k\in\N,\\
  0\pmod4,&\text{otherwise}.
 \end{cases}
\]
\end{theorem}

\begin{proof}
The construction of $L$ works uniquely over any commutative ring.
It therefore commutes with reduction, as do $R=x-2xL$ and its genuine
compositional inverse. Modulo two, equation \eqref{eq:squareL} becomes
$L+L^2=x$. Lemma~\ref{lem:binary} identifies its support with the powers
of two.

Now work modulo four. Since $R=x-2xL\equiv x\pmod2$, substitution gives
$L(R)\equiv L\pmod2$, and Lemma~\ref{lem:double} gives
$2L(R)=2L$. Consequently
\begin{align*}
 R(R)&=R-2R\,L(R)=R-2R\,L\\
     &=(x-2xL)(1-2L)=x-4xL+4xL^2=x.
\end{align*}
Thus $R$ is its own compositional inverse modulo four.
Uniqueness of compositional inverses identifies the reduction of $A$
with $R$:
\[
 A\equiv x-2xL\pmod4.
\]
For $n>1$, the coefficient is $-2\ell_{n-1}$ modulo four. This equals
two precisely when $\ell_{n-1}$ is odd, which holds exactly when
$n-1=2^k$; otherwise it is zero.
\end{proof}

The result is the single modulus-four conjecture in A389542
\cite{oeis542}. The entire involution argument uses multiplication by
two and reduction to characteristic two; it makes no cancellation of
two in $\Z/4\Z$.


\section{Separate Lean proof and checked scope}

The entry file is \nolinkurl{lean/Entries/A389542.lean}. Its explicit
theorem statements and proof terms use the namespace
\nolinkurl{FinalCollection.A389542}. The complete imported implementations
and their shared mathematical dependencies are included in the collection.
The following local theorem names give the proof map:

\begin{multicols}{2}\raggedright\begin{itemize}\setlength{\itemsep}{2pt}\setlength{\parsep}{0pt}
\item \nolinkurl{auxiliary_exists_unique}.
\item \nolinkurl{normalized}.
\item \nolinkurl{auxiliary_definition}.
\item \nolinkurl{inverse_equation}.
\item \nolinkurl{inverse_right}.
\item \nolinkurl{inverse_left}.
\item \nolinkurl{source_equation}.
\item \nolinkurl{source_exists_unique}.
\item \nolinkurl{source_unique}.
\item \nolinkurl{auxiliary_parity}.
\item \nolinkurl{inverse_involution}.
\item \nolinkurl{inverse_mod_four}.
\item \nolinkurl{series_mod_four}.
\item \nolinkurl{residue_mod_four}.
\item \nolinkurl{exceptional_iff}.
\item \nolinkurl{four_dvd_iff}.
\end{itemize}\end{multicols}

The main implementation module is \nolinkurl{Round3.SquareReversion}.
The catalogue records the full original endpoint names and the new entry
theorems. The project pins Lean 4.19.0 and Mathlib revision
\nolinkurl{c44e0c8ee63ca166450922a373c7409c5d26b00b}. The collection
report records source recompilation and transitive axiom checks. Accepted
axioms are propositional extensionality, classical choice, and quotient
soundness. All index conclusions above are unbounded; exact finite
coefficient checks in the package provide separate corroboration.

\paragraph{Limits and status.}
Source uniqueness is proved over the integers, not over arbitrary rings with two-torsion. No analytic value or radius statement is claimed.
The manuscript consolidates completed results; it does not assert global
novelty, exclusive priority, peer review, OEIS acceptance, or a new submission.
Standard algebraic ingredients are not claimed as new discoveries.

\begin{thebibliography}{9}
\raggedright
\bibitem{oeis542}
Paul D. Hanna, \emph{OEIS A389542}, 1 November 2025.
\url{https://oeis.org/A389542}. Retrieved 24 October 2026.
\end{thebibliography}
\end{document}
